Our experiments demonstrate significant improvements across all measured dimensions when using the full Discord OpenClaw framework compared to baselines.

In the User Engagement Study (Table \ref{tab:engagement}), proactive behavior reduced response latency by 33.3\% (4.2s to 2.8s, p=0.002), increased engagement rate by 77.1\% (0.35 to 0.62, p=0.001), extended session length by 67.1\% (8.5min to 14.2min, p=0.003), and improved user satisfaction by 28.1\% (3.2 to 4.1, p=0.001). These results strongly support our first claim that proactive behavior enhances user engagement in communication platforms.

The Safety and Trust Evaluation (Table \ref{tab:safety}) showed that requiring approval for tool execution increased safety ratings by 53.6\% (2.8 to 4.3, p=0.0001), trust scores by 41.9\% (3.1 to 4.4, p=0.0001), and prevented harmful actions by 1300\% (0.2 to 2.8, p=0.0001), while only increasing task completion time by 30.8\% (5.2s to 6.8s, p=0.001) and having no significant impact on user satisfaction (p=0.08). This confirms our second claim that approval workflows enhance safety and trust without significantly reducing utility.

The External Knowledge Integration Study (Table \ref{tab:knowledge}) revealed that NotebookLM integration improved factual accuracy by 30.9\% (0.68 to 0.89, p=0.0001), reduced hallucination rates by 68.2\% (0.22 to 0.07, p=0.0001), improved relevance scores by 9.3\% (0.75 to 0.82, p=0.001), improved citation quality from 0.00 to 0.76 (p=0.0001), and increased user satisfaction by 19.4\% (3.6 to 4.3, p=0.0001). These results validate our third claim that external reasoning integration improves answer quality and reduces hallucinations.

Finally, the Component Ablation Study (Table \ref{tab:ablation}) demonstrated that all components contribute significantly to overall performance. The full system achieved an average performance score of 8.2/10. Removing proactive behavior decreased performance to 6.1/10 (25.6\% reduction), removing approval workflows decreased performance to 5.8/10 (29.3\% reduction), removing external reasoning integration decreased performance to 6.5/10 (20.7\% reduction), and the reactive-only baseline scored 4.2/10 (48.8\% reduction from full system). All differences were statistically significant (p=0.0001 for ANOVA tests). This shows that each component makes a unique contribution, with proactive behavior and approval workflows having the largest individual impacts.